% BHU PYQ -- Transcribed & Corrected
% Paper: CIMSE21: Computational Mathematics with SageMath
% Exam: B.Sc. (Hons.) Semester II Examination 2024-25 (NEP)
% Ref Code: CE/CONF./N/35
% Category: SEC (DST-CIMS)
% Department: Computational Mathematics / Mathematics

\documentclass[11pt,a4paper]{article}
\usepackage[a4paper,top=2.5cm,bottom=2.5cm,left=2.8cm,right=2.8cm,headheight=14pt]{geometry}
\usepackage{amsmath,amssymb}
\usepackage[T1]{fontenc}
\usepackage[utf8]{inputenc}
\usepackage{lmodern,microtype,fancyhdr,enumitem,hyperref,lastpage,parskip,listings}

\hypersetup{
    colorlinks=true,
    urlcolor=black,
    linkcolor=black,
    pdftitle={CIMSE21: Computational Mathematics with SageMath -- BHU Sem II 2024-25 (NEP SEC)}
}

\pagestyle{fancy}
\fancyhf{}
\renewcommand{\headrulewidth}{0.4pt}
\renewcommand{\footrulewidth}{0.4pt}
\fancyhead[L]{\small \textit{Banaras Hindu University}}
\fancyhead[R]{\small \textit{CIMSE21: SageMath (2024-25)}}
\fancyfoot[C]{\small Page \thepage\ of \pageref{LastPage}}

\newcommand{\pts}[1]{\hfill \textit{[#1]}}

\begin{document}

\begin{center}
  {\large\bfseries BANARAS HINDU UNIVERSITY}\\[2pt]
  {\normalsize B.Sc. (Hons.) Semester II Examination 2024-25 (NEP)}\\[6pt]
  \rule{0.8\linewidth}{0.5pt}\\[6pt]
  {\large\bfseries SEC (DST-CIMS)}\\[4pt]
  {\large\bfseries Paper Code: CIMSE21 : Computational Mathematics with SageMath}\\[2pt]
  {\small Ref. No.: CE/CONF./N/35}\\[4pt]
  {\normalsize Time: 2$\frac{1}{2}$ Hours\hfill Maximum Marks: 70}
\end{center}

\rule{\linewidth}{0.4pt}

\smallskip

\noindent \textit{Instructions: Solve any five questions including Question 1 which is compulsory. Each question carries equal marks.}

\bigskip

\noindent \textbf{Question 1.} Write the codes in SageMath for the following mathematical operations and then solve the corresponding question without using codes: \pts{14}
\begin{enumerate}[label=(\alph*)]
  \item To find the rank of a matrix and hence find the rank of the matrix $A = \begin{pmatrix} 4 & 8 & 16 \\ 1 & 2 & 4 \\ 0 & 0 & 1 \end{pmatrix}$.
  \item To find the adjoint of a matrix A and hence find the adjoint of the matrix $A = \begin{pmatrix} 5 & 3 \\ 3 & 2 \end{pmatrix}$.
  \item To find the remainder when an integer $a$ is divided by a positive integer $b$ and hence find the remainder when $(2025)^3$ is divided by 47.
  \item To solve a system of linear equations $2x + y = 6, 5x - y = 4$ in two variables and hence solve it manually.
  \item To plot the curve $y^2 = x + 2y$ in $\mathbb{R}^2$ and hence plot it.
  \item To find the dot product of two vectors in $\mathbb{R}^3$ and hence find the dot product of two vectors which represent the sides of an equilateral triangle of side length 1.
  \item To plot a vector field in $\mathbb{R}^2$ and hence plot the vector field $V(x, y) = (-y, x)$ in the region $\{(x, y) \in \mathbb{R}^2 \mid x^2 + y^2 < 1\}$.
\end{enumerate}

\bigskip

\noindent \textbf{Question 2.}
\begin{enumerate}[label=(\alph*)]
  \item Write a 3 dimensional Cartesian equation of a line in $xz$ plane, passing through origin and making equal angle with the two axes. Write the codes in SageMath to plot this line. \pts{7}
  \item What is rref of a matrix? Write a SageMath program to find the rref of the matrix $A = \begin{pmatrix} 2 & 3 & 4 \\ 3 & 4 & 5 \\ 4 & 5 & 6 \end{pmatrix}$. Compute manually to find its output. \pts{7}
\end{enumerate}

\bigskip

\noindent \textbf{Question 3.}
\begin{enumerate}[label=(\alph*)]
  \item State Lagrange mean value theorem and verify it for the function $f(x) = x^2 + 3x + 2$ on the interval $[-1, 1]$. Can you interpret it geometrically? Write a SageMath program to check the validity of your computation. \pts{7}
  \item Write a SageMath program to plot the tangent line to a circle of your choice at any of its point so that the colours of the circle and the tangent line appear differently. \pts{7}
\end{enumerate}

\bigskip

\noindent \textbf{Question 4.}
\begin{enumerate}[label=(\alph*)]
  \item Define eigen value and eigen vector of a square matrix A and hence find them for the matrix $A = \begin{pmatrix} 2 & 3 \\ 2 & 1 \end{pmatrix}$. Write codes in SageMath to find eigen values and eigen vectors of a square matrix A. \pts{7}
  \item Define diagonalizability of a square matrix. Determine whether the matrix A, as defined above, is diagonalizable? Write SageMath program to check the diagonalizability of a square matrix. \pts{7}
\end{enumerate}

\bigskip

\noindent \textbf{Question 5.}
\begin{enumerate}[label=(\alph*)]
  \item Find the area of the region enclosed by the parabola $y^2 = 4x$ and the line $y = 2x$ and write the SageMath code to compute it. \pts{7}
  \item What will be the output of the following SageMath program: \pts{7}\\
  \texttt{integrate((1/(1 + e\textasciicircum(sin(x)))), (x, 0, 2*pi))}
\end{enumerate}

\bigskip

\noindent \textbf{Question 6.}
\begin{enumerate}[label=(\alph*)]
  \item Find the differential equation of all circles which touches the line $y = x$ at the origin and hence by solving it find the equation of that circle which has intercept 1 on positive x-axis. Also write SageMath program to solve the differential equation. \pts{7}
  \item Find the differential equation of the curves passing through the origin given that the slope of the tangent to the curve at any point $(x, y)$ is equal to the sum of the coordinates of the point and hence write SageMath code to solve it. \pts{7}
\end{enumerate}

\bigskip

\noindent \textbf{Question 7.}
\begin{enumerate}[label=(\alph*)]
  \item Find the volume of the parallelopiped formed by vectors $(1,1,1)$, $(1,1,0)$ and $(2,3,4)$ in $\mathbb{R}^3$ as its adjacent edges. Also write the codes in SageMath to find the volume of the parallelopiped, whose three adjacent edges are given as vectors in $\mathbb{R}^3$. \pts{7}
  \item Write SageMath program to plot a right circular cylinder with z-axis as its axis and radius 1 in two different ways. \pts{7}
\end{enumerate}

\bigskip
\begin{center}
  \textbf{***}
\end{center}

\end{document}
